
\documentclass[12pt]{article}
\usepackage[margin=1in]{geometry}
\usepackage{amsmath,amssymb,booktabs,longtable}
\usepackage{hyperref}
\usepackage{multirow}
\usepackage[margin=1in]{geometry}
\usepackage[T1]{fontenc}
\usepackage[utf8]{inputenc}
\usepackage{lmodern}
\usepackage{microtype}
\usepackage{enumitem}
\usepackage[hidelinks]{hyperref}
\setlength{\parindent}{0pt}
\setlength{\parskip}{0.7em}
\setlist{nosep}
\title{Lecture Summary and DNA Structure}
\author{Zachary Tullier}
\date{7/28/26}
\begin{document}
\maketitle

\section*{Q1. }

This lecture explains how radiation interacts with matter, how those
interactions affect biological molecules and cells, and how the same
principles are applied in medicine, environmental science, and space travel.

\section{Radiation and Its Sources}

Radiation spans a broad electromagnetic spectrum. Low-frequency radio waves,
microwaves, infrared radiation, and visible light generally do not carry enough
energy to ionize matter. X-rays and gamma rays have much shorter wavelengths
and greater photon energies, allowing them to remove electrons from atoms and
molecules.

Radiation comes from both natural and artificial sources:

\begin{itemize}
    \item The Sun produces visible light, ultraviolet radiation, solar
    neutrinos, gamma rays, protons, and alpha particles.
    \item High-energy cosmic rays originate outside Earth and produce showers
    of secondary particles when they collide with atmospheric nuclei.
    \item Radioactive decay, nuclear reactions, particle accelerators, medical
    equipment, and nuclear technology provide terrestrial or artificial
    sources.
\end{itemize}

\section{Ionizing Versus Non-Ionizing Radiation}


\begin{tabular}{ |p{3cm}||p{3cm}|p{3cm}|p{3cm}|  }
 \hline
 \multicolumn{4}{|c|}{Country List} \\
 \hline
 \textbf{Possible biological effects} & \textbf{Category} & \textbf{Examples}& \textbf{Primary interaction} \\
 \hline
 Non-ionizing &
Radio waves, microwaves, infrared, visible light, and most ultraviolet
radiation &
Excites molecules or produces heat without normally removing electrons &
Thermal effects, photochemical reactions, and some forms of DNA damage \\
 Ionizing &
High-energy ultraviolet radiation, X-rays, gamma rays, alpha and beta
particles, protons, and neutrons &
Removes electrons and produces ions and reactive chemical species &
DNA breaks, mutations, cell death, tissue injury, and cancer \\
 
 \hline
\end{tabular}


Although UVA and UVB are usually classified as non-ionizing, they remain
biologically harmful. UVA efficiently generates reactive oxygen species,
whereas UVB can produce direct photochemical DNA lesions. These processes
oxidize bases such as guanine and form thymine dimers and other cross-links.
Such structural distortions can interfere with DNA replication and
transcription.

\section{Radioactive Decay and Radiation Measurement}

The lecture reviewed several nuclear transformations:

\begin{itemize}
    \item Alpha decay releases a helium nucleus.
    \item Beta-minus decay releases an electron and converts a neutron into a
    proton.
    \item Positron emission releases a positively charged electron.
    \item Electron capture occurs when a nucleus absorbs an orbital electron.
    \item Gamma decay releases electromagnetic energy without changing the
    element's atomic number.
\end{itemize}

Different units describe different aspects of radiation:

\section*{Q2. DNA Structure and Rectification Process}


\textit{Note: This document interprets ``rectification'' as DNA repair. DNA
replication is a separate process.}

\section{Structure of DNA}

DNA, or deoxyribonucleic acid, is composed of two antiparallel polynucleotide
strands twisted into a double helix.

Each DNA nucleotide contains:

\begin{itemize}
    \item a deoxyribose sugar;
    \item a phosphate group; and
    \item one nitrogenous base: adenine (A), thymine (T), guanine (G), or
    cytosine (C).
\end{itemize}

The sugar and phosphate groups form the outer backbone of each strand. The
nitrogenous bases face inward and pair through hydrogen bonds:

\begin{itemize}
    \item adenine pairs with thymine through two hydrogen bonds; and
    \item guanine pairs with cytosine through three hydrogen bonds.
\end{itemize}

The two strands run in opposite directions. One strand runs from
\(5' \rightarrow 3'\), while the other runs from \(3' \rightarrow 5'\).
Complementary base pairing allows each strand to serve as a template during
DNA replication and repair.

Within eukaryotic cells, DNA is wrapped around histone proteins to form
nucleosomes. Nucleosomes are organized into chromatin, which can be further
condensed into chromosomes.

\section{DNA Repair Process}

DNA is continually damaged by replication errors, ultraviolet light, ionizing
radiation, reactive oxygen species, and chemical agents. Repair generally
involves three stages:

\begin{enumerate}
    \item The cell recognizes the damaged or incorrectly paired DNA.
    \item The damaged base or section of DNA is removed or corrected.
    \item DNA polymerase restores the correct sequence, and DNA ligase seals
    the sugar--phosphate backbone.
\end{enumerate}

\subsection{Direct Reversal}

Direct reversal repairs certain chemical modifications without removing the
affected nucleotide. For example, the enzyme
\textit{O}\textsuperscript{6}-methylguanine-DNA methyltransferase (MGMT) can
remove particular alkyl groups from guanine.

\subsection{Base-Excision Repair}

Base-excision repair corrects small, chemically altered bases:

\begin{enumerate}
    \item A DNA glycosylase recognizes and removes the damaged base.
    \item An endonuclease cuts the DNA backbone at the resulting site.
    \item DNA polymerase inserts the correct nucleotide.
    \item DNA ligase seals the repaired strand.
\end{enumerate}

\subsection{Nucleotide-Excision Repair}

Nucleotide-excision repair removes bulky lesions that distort the DNA double
helix, including ultraviolet-induced pyrimidine dimers. The cell excises a
short section of the damaged strand, DNA polymerase synthesizes a replacement
section, and DNA ligase seals the remaining break.

\subsection{Mismatch Repair}
Mismatch repair corrects incorrectly paired bases and small insertion or
deletion errors that escape proofreading during DNA replication. The repair
system identifies the newly synthesized strand, removes the section containing
the error, and replaces it with the correct sequence.

\subsection{Single-Strand-Break Repair}

When only one DNA strand is broken, repair proteins process the damaged ends,
DNA polymerase replaces any missing nucleotides, and DNA ligase reconnects the
strand. The intact complementary strand helps preserve the correct genetic
information.

\subsection{Double-Strand-Break Repair}

Double-strand breaks are especially dangerous because both DNA backbones are
severed. Cells use two principal repair pathways:

\begin{itemize}
    \item \textbf{Homologous recombination} uses an undamaged homologous DNA
    sequence, usually the sister chromatid, as a template. This process is
    generally highly accurate.
    \item \textbf{Non-homologous end joining} reconnects the broken ends
    directly. It is faster and does not require a homologous template, but it
    may introduce small mutations at the repair site.
\end{itemize}

\section{Cellular Response to Severe Damage}

Cell-cycle checkpoints temporarily stop cell division while repair occurs. If
the damage cannot be repaired, the cell may enter permanent growth arrest or
undergo programmed cell death, known as apoptosis. If defective repair allows
a damaged cell to survive and continue dividing, mutations may accumulate and
contribute to cancer.

In radiation biophysics, double-strand breaks and clustered lesions are
particularly important because they are more difficult to repair accurately
than isolated base damage or single-strand breaks.

\section{Conclusion}

DNA's complementary double-stranded structure provides both stable storage of
genetic information and a template for error correction. Multiple repair
pathways protect this information by recognizing different forms of damage,
restoring the correct nucleotide sequence, and resealing the DNA backbone.
The success or failure of these pathways strongly influences mutation,
cellular survival, aging, and disease development.

\end{document}
